\section{Linéariser le produit de fonction trigonométriques}

On utilise les formules $\sin(a)\sin(b)=\frac{1}{2}\left(\cos(a-b)-\cos(a+b)\right)$, $\cos(a)\cos(b)=\frac{1}{2}\left(\cos(a-b)+cos(a+b)\right)$ et $\sin(a)\cos(b)=\frac{1}{2}\left(\sin(a-b)+\sin(a+b)\right)$.

\section{Déterminer le spectre de produits et sommes de sinus}
\begin{enumerate}
    \item S'il y a des produits de fonctions trigonométriques dépendant du temps, les linéariser.
    \item Pour chaque terme, représenter un pic à la fréquence et l'amplitude correspondante sur le spectre.
\end{enumerate}

